\documentclass[UTF8,fontset=none,11pt,a4paper]{ctexart}
\usepackage{ipara-handbook}
\handbooksetup
  {DS-02 栈、队列与受限访问}
  {408 · 数据结构 DS-02}
  {Restriction · Order · State · Invariant}
\tikzset{
  state/.style={draw=iparaDeep,minimum width=2.6cm,minimum height=0.95cm,align=center,fill=white,font=\small},
  access/.style={-{Latex[length=2mm]},thick,draw=iparaDeep},
  forbid/.style={-{Latex[length=2mm]},dashed,draw=iparaRule}
}
\begin{document}
\handbooktitle
  {栈、队列与受限访问}
  {减少端点自由度，怎样生成可推理的时间秩序}
  {408 · 数据结构 Topic DS-02}

\begin{corebox}{母问题：为什么“少允许一些操作”反而能表达过程？}
一般线性表允许在许多位置访问和修改；栈、队列与双端队列主动收紧这个自由度。它们的生成链是
\[
\begin{aligned}
\text{Access Restriction}&\to\text{Temporal Order}\to\text{State Transition}\\
&\to\text{Invariant}\to\text{Application Contract}.
\end{aligned}
\]
这里的箭头依次表示因果生成和检查顺序：先规定哪些端点可进、可出，才能推出元素被取出的时间秩序；每次操作再更新状态字段，并保持空、满、首尾指针和元素次序不变量。应用只依赖这份契约，不应依赖底层是数组还是链表。
\end{corebox}

\begin{methodbox}{本册定位与 Owner 边界}
\begin{itemize}
\item \kw{Owns：}栈、队列、循环队列、链队列、双端队列的 ADT 契约、物理状态、操作生命周期、空满判据与复杂度。
\item \kw{Uses：}DS01 的顺序/链接表示；Atlas 的对象—操作—不变量—成本语言。
\item \kw{Outputs：}向 DS-B01 输出可替换的 Frontier 容器契约；向 DS04/DS07 提供遍历所需的取出顺序。
\item \kw{Stops：}递归语义、表达式文法、树/图遍历正确性和单调结构的淘汰规则不由本册拥有；这里只说明它们为何能够调用受限访问容器。
\end{itemize}
\end{methodbox}

\tableofcontents\newpage

\section{从线性表到受限访问 ADT}
若线性表的状态是一个序列 $S=(a_1,\ldots,a_n)$，那么“栈”和“队列”的区别不在元素类型，而在允许观察和删除哪一端：
\begin{center}\small
\begin{tabularx}{\linewidth}{L{2.2cm}L{2.4cm}L{2.4cm}YY}
\toprule
ADT & 允许进入 & 允许离开/观察 & 推出的次序 & 典型过程含义\\
\midrule
栈 Stack & 顶端 & 顶端 & LIFO，后进先出 & 最近未完成的上下文\\
队列 Queue & 队尾 & 队头 & FIFO，先进先出 & 最早等待的任务\\
双端队列 Deque & 两端 & 两端 & 由调用者选择 & 可同时模拟栈和队列\\
\bottomrule
\end{tabularx}
\end{center}

限制不是性能结论。栈的 \code{push}/\code{pop} 可以由顺序表尾部实现，也可以由链表头部实现；只要返回序列一致，二者兑现的是同一 ADT。表示改变的是容量、分配、局部性和维护成本，不是 LIFO/FIFO 的抽象结果。

\begin{examplebox}{完整生成例：三项任务按栈和队列退出}
依次加入 $A,B,C$。若只允许从加入端退出，则状态轨迹为
\[
[]\to[A]\to[A,B]\to[A,B,C]\to[A,B]\to[A]\to[],
\]
退出序列是 $C,B,A$。若一端加入、另一端退出，则退出序列是 $A,B,C$。同一组元素没有改变，访问限制却确定了时间秩序；这正是应用选择容器时真正购买的语义。
\end{examplebox}

\section{栈：顶端是唯一活动边界}
把栈状态写作 $(data,size,capacity)$ 或 $(top,size)$。顺序栈的合法状态是
\[
0\le size\le capacity,\qquad
\operatorname{top}=data[size-1]\quad(size>0).
\]
执行 \code{push(x)} 时先判满，再写入 $data[size]$ 并递增长度；执行 \code{pop()} 时先判空，再递减长度并读取原顶元素。\kw{先减后读}与\kw{先读后减}都可正确实现，但下标和返回值必须指向同一个旧顶元素。

链栈把头节点作为栈顶：压栈令新节点的后继指向旧顶，再更新顶指针；出栈先保存旧顶，再跨过它并释放。若把尾节点当栈顶而底层又是单链表，出栈还要寻找前驱，便丢掉了预期的 $O(1)$ 局部更新。

\begin{center}
\begin{tikzpicture}[node distance=0.7cm]
\node[state](contract){栈契约\\只访问 top};
\node[state,right=of contract](array){顺序栈\\size 是下一空位};
\node[state,right=of array](linked){链栈\\头指针是 top};
\draw[access] (contract)--node[above,font=\scriptsize]{表示 A}(array);
\draw[access] (contract)--node[above,font=\scriptsize]{表示 B}(linked);
\end{tikzpicture}
\end{center}

\section{队列：两端分工生成 FIFO}
队列允许在队尾 \code{enqueue}、在队头 \code{dequeue}。链队列通常保存 $(front,rear,size)$：入队接到 $rear$ 后，出队从 $front$ 删除。最危险的不是一般节点，而是最后一个节点出队：新 $front$ 为空时，$rear$ 也必须复位为空，否则“空队列”会残留悬空尾指针。

\subsection{为什么顺序队列需要循环化}
若数组队列只让 $front,rear$ 单调右移，前部空槽无法复用，会出现“数组还有空间但队尾到界”的假溢出。循环队列把下标放在模 $M$ 的环上：
\[
\operatorname{next}(i)=(i+1)\bmod M.
\]
这不是为了背取模公式，而是把已出队的槽重新纳入可用状态空间。

\subsection{空与满必须增加一份信息}
只保存 $front,rear$ 且二者都指向下一操作位置时，$front=rear$ 既可能表示空，也可能表示绕一圈后的满。实现必须选择一种消歧协议：
\begin{center}\small
\begin{tabularx}{\linewidth}{L{2.4cm}YYY}
\toprule
协议 & 空判据 & 满判据 & 代价/适用点\\
\midrule
牺牲一个槽 & $front=rear$ & $next(rear)=front$ & 可用容量为 $M-1$\\
长度计数 & $size=0$ & $size=M$ & 多维护一个计数字段\\
标志/最近操作 & $front=rear$ 且标志为空 & $front=rear$ 且标志为满 & 分支语义更复杂\\
\bottomrule
\end{tabularx}
\end{center}
本册伴随实现选择长度计数，因此合法状态还要求
\[
0\le size\le M,\qquad rear\equiv front+size\pmod M.
\]
不能把该实现的 $size=M$ 满判据与“牺牲一个槽”实现的 $next(rear)=front$ 混用。

\section{双端队列：两端自由，但状态仍受约束}
双端队列允许 \code{push\_front}/\code{push\_back} 与 \code{pop\_front}/\code{pop\_back}。栈和普通队列都可由它限制操作集合得到：
\[
\text{Deque}\xrightarrow{\text{只在同端进出}}\text{Stack},\qquad
\text{Deque}\xrightarrow{\text{一端进、另一端出}}\text{Queue}.
\]
这说明 Deque 的 ADT 能力更强，却不等于所有应用都应使用它。更宽的接口允许更多非法调用组合；若应用只需要 FIFO，直接暴露 Queue 契约更能保护时间秩序。

循环双端队列仍维护 $rear\equiv front+size\pmod M$。前端插入要先令 $front$ 向前回绕再写入；后端删除要先令 $rear$ 向前回绕再读取。顺序写反会读到“下一空位”，产生典型的 off-by-one 错误。

\section{操作序列、共享空间与表达式：受限访问的题面变体}
\subsection{合法出栈序列不是“看起来像后进先出”}
给定输入序列 $1,2,\ldots,n$，每一步可以继续入栈，也可以在栈非空时出栈。判定一个候选输出序列时，维护“已经读入但尚未输出”的栈：
\begin{enumerate}
\item 若栈顶等于当前候选输出，立即出栈并前进输出指针；
\item 否则若输入尚未耗尽，继续把下一个输入压栈；
\item 若输入耗尽且栈顶仍不是候选输出，序列非法。
\end{enumerate}
因此合法性是一个状态可达性问题，不是只检查“有没有重复”。例如 $1,2,3$ 可以输出 $2,3,1$（先压 $1,2$，出 $2$，再压 $3$，出 $3,1$），但不能输出 $3,1,2$，因为输出 $3$ 时 $1,2$ 都在它下面，随后不可能先取出 $1$。计数题还必须区分“输入次序固定的栈排列”与“任意元素可直接访问的排列”，后者不再是栈问题。

双端队列题要先声明输入限制：输入只能从一端进入、输出只能从一端离开，还是两端都可用。不同限制会产生不同可行序列；把“输入受限 Deque”和普通 Queue 混成同一张表，会丢掉题目真正考察的端点约束。

\subsection{共享栈：两个增长方向共享一个容量不变量}
在同一数组中让两个栈相向增长，典型状态为 $top_1$ 指向左栈顶、$top_2$ 指向右栈顶。若两端均把“下一空位”作为边界，栈满条件是
\[
top_1+1=top_2.
\]
这里的“满”不是某一个栈达到数组末端，而是两条有效区域即将相撞。左栈压入会增大 $top_1$，右栈压入会减小 $top_2$；出栈只改变对应一端。共享栈的收益是两个栈可以互相借用暂时空出的槽位，代价是每个操作都必须检查全局相撞条件。若题目采用“top 指向栈顶元素”而不是“下一空位”，公式会整体平移一格，必须先画空栈状态再推导，不能直接背上式。

\subsection{中缀、后缀与表达式求值：运算符栈维护的是未决语义}
表达式题的核心不是把字符串换一种写法，而是把“尚未决定结合次序的运算符”保存到栈中。中缀转后缀时维护两个区域：输出序列已经确定的 token，以及运算符栈中等待比较优先级的 token。
\begin{enumerate}
\item 扫到操作数，直接输出；
\item 扫到左括号，压栈并把它视为新的优先级边界；
\item 扫到右括号，不断弹出直到左括号，左右括号本身不输出；
\item 扫到运算符 $op$，若栈顶运算符优先级不低于 $op$ 且不跨越左括号，就弹出栈顶，再压入 $op$；
\item 扫描结束后把栈中剩余运算符依次弹出。
\end{enumerate}
“大于等于”而非“大于”决定同优先级运算符的左结合；若题目允许右结合的幂运算，比较规则必须单独写出。表达式求值时，操作数栈每遇到二元运算符弹出右操作数 $b$、再弹出左操作数 $a$，计算 $a\;op\;b$ 后压回；弹出顺序不能交换。除法、负号、括号和非法 token 都是边界分支，不应被“栈算法”四个字掩盖。

\subsection{递归调用栈：每一帧保存尚未完成的未来}
递归不是“函数自己调用自己”的语法现象，而是把每次调用的局部状态压入系统栈。每个递归函数至少要回答三个问题：
\begin{enumerate}
\item 参数表示的子问题是什么，规模怎样变小；
\item 终止条件是否覆盖空结构、叶节点和目标已命中；
\item 子调用返回后，当前帧还要执行什么合并、回溯或访问动作。
\end{enumerate}
对树遍历，访问语句放在递归调用前、两次调用之间或两次调用之后，分别生成前序、中序和后序；这正是“返回时机”与栈帧恢复顺序的关系。把递归改写为显式栈时，不能只保存节点，还要保存“下一步要执行哪个子调用”的程序计数状态，否则只能得到某一种遍历而不能等价模拟原递归。

递归空间复杂度应计入最大同时存在的调用帧数。递归深度为 $h$ 的树遍历使用 $O(h)$ 栈空间；若树退化成链，$h=n$，不能仅凭“每个节点访问一次”写成 $O(1)$ 额外空间。尾递归是否能由编译器消除是实现扩展，不能替换 408 题目要求的抽象调用栈分析。

\subsection{括号匹配与单调结构的边界}
括号匹配是最小的栈应用：遇到左括号压入类型，遇到右括号时必须检查栈非空且栈顶类型匹配，扫描结束后栈必须为空。三类失败分别是“右括号无对应左括号”“类型不匹配”“仍有未闭合左括号”。

单调栈、单调队列是在普通受限访问之外增加了候选淘汰不变量：新元素到来时，所有不可能再成为答案的旧元素被从端点弹出。它们仍使用 DS02 的端点操作，但“为什么可以淘汰”属于具体问题的支配关系/窗口约束，不属于普通栈或队列定义。规则进入考试控制或对应算法题，不在本册扩成第四种容器。

\subsection{顺序栈的 top 约定：先画空栈，再写代码}
顺序栈常见两种等价口径：\code{top} 指向当前栈顶元素，空栈为 $top=-1$；或 \code{top} 指向下一空位，空栈为 $top=0$。前者入栈是“先加后写”、出栈是“先读后减”；后者入栈是“先写后加”、出栈是“先减后读”。两种口径都能得到同一 LIFO 结果，错误只发生在初始化、判空、判满和读写顺序混用时。动态顺序栈还要明确 $base$ 指向连续存储区、$top-base$ 是元素数量；链栈则没有固定上界，失败条件是分配失败而不是数组上溢。

\subsection{合法出栈序列：模拟与禁用模式}
输入顺序固定为 $1,2,\ldots,n$ 时，出栈过程可看作“从输入流继续压入”与“从栈顶输出”两种动作的交错。最可靠的判定算法是维护输入指针、辅助栈和输出指针：栈顶等于目标就弹出；否则继续压入尚未读取的输入；输入耗尽且栈顶仍不匹配即失败。任何候选输出都必须同时满足每个元素只出现一次、输入次序不被跳过以及每次输出确实在栈顶。

还有一个可用于复核的禁用模式：若输出序列中存在 $a<b<c$，且 $a$ 在 $b$、$c$ 之前而 $c$ 在 $b$ 之前（相对次序呈 $a,c,b$），则它不可能由固定输入栈产生；$a$ 压入后 $b$ 在它上面，若先取出更晚的 $c$，$b$ 会被挡住而不能随后取出。计数时，固定输入顺序的合法出栈序列数是 Catalan 数
\[
C_n=\frac{1}{n+1}\binom{2n}{n},
\]
它等价于“入栈/出栈前缀中任意时刻出栈数不超过入栈数”的 Dyck 路径计数；不要把它与任意排列的 $n!$ 混淆。括号匹配、二叉树形状和合法栈操作序列中出现同一递推，是因为它们都在维护一个不能先耗尽的未决边界。

\subsection{双端队列的三种权限与输出序列}
题面中的“受限双端队列”必须先标出输入端和输出端：
\begin{center}\small
\begin{tabularx}{\linewidth}{L{2.1cm}YY}
\toprule
类型 & 允许操作 & 判题重点\\\midrule
输入受限 & 只能从一端输入，两端可输出 & 每个新元素进入固定端，输出时可选两端\\
输出受限 & 两端可输入，只能从一端输出 & 需模拟两端输入如何排队到唯一出口\\
完全双端 & 两端均可输入/输出 & 能力最强，但若题面未限制则不能擅自增加操作\\
\bottomrule
\end{tabularx}\end{center}
输入受限时，候选输出序列可用首尾剥离规则快速复核：读入下一个元素前，目标值若不是当前两端之一就必须继续读；若两端均不匹配且输入耗尽则失败。输出受限时没有同样的一眼规则，按“从左端/右端输入”逐步模拟更稳妥。栈和队列只是双端队列的两个操作子集，不能因为某序列能由 Deque 产生就断言它能由 Stack 或 Queue 产生。

\subsection{循环队列三种消歧协议与元素个数}
同一数组环可以有三套合法设计，必须整套使用：
\begin{enumerate}
\item 牺牲一个槽：$front=rear$ 表示空，$next(rear)=front$ 表示满，最多存 $M-1$ 个元素；
\item 保存 $size$：空为 $size=0$，满为 $size=M$，且 $size=(rear-front+M)\bmod M$（若 $rear$ 指下一空位）；
\item 保存 $tag$：入队/出队后端点相遇时分别把标志置为“满/空”，容量可用满 $M$ 个槽，但状态转移多一条标志维护。
\end{enumerate}
若 $rear$ 指向队尾元素而不是下一空位，元素个数公式和入/出队顺序会整体平移；正确做法是先画“空队列”和“恰好一个元素”的指针图，再推公式。所谓“假溢出”只属于非循环顺序队列前端空槽未复用的实现缺陷，不属于逻辑队列容量本身。

\subsection{表达式的完整栈模型：优先级、结合性与多栈}
中缀转后缀用运算符栈时，栈内优先级与栈外优先级需要区分：左括号入栈后形成新的边界，右括号触发弹栈直到左括号；同优先级左结合运算符遇到新运算符时应先弹出，右结合运算符（若题目允许）则保留栈顶。扫描结束必须把剩余运算符弹出，否则后缀表达式缺少尾部运算。

后缀求值只需操作数栈：遇操作数压入，遇二元运算符先弹右操作数 $b$、再弹左操作数 $a$，计算 $a\,op\,b$。中缀直接求值则维护“操作数栈 + 运算符栈”两个状态；新运算符优先级不高于栈顶时先完成栈顶运算。前缀求值从右向左扫描，弹出顺序仍是先得到左/右操作数后按运算符语义组合。表达式树把三种形式统一为前序/中序/后序遍历；中缀输出必须按子树优先级补括号，不能只输出一次中序遍历就声称无歧义。

\subsection{递归栈与显式栈的等价条件}
一次函数调用要在系统栈帧中保存返回地址、参数、局部变量和下一条待执行位置；递归返回时恢复该帧，所以“访问语句放在递归调用前/中/后”分别对应前序/中序/后序。把递归改成显式栈时，若原函数在返回后还要处理右子树，栈项必须额外保存节点和阶段标记（例如“左子树已完成，下一步处理右子树”），只压入节点指针会丢失控制状态。递归时间看调用树节点总数，递归空间看同时存在的根到当前节点路径；两者可以分别是 $O(n)$ 和 $O(h)$。

\section{代码把访问契约落实为状态机}
完整 C++17 实现位于 \codepath{code/ds02_restricted_access.hpp}，测试位于 \codepath{code/ds02_restricted_access_test.cpp}。以下代码块直接抽取真实源文件，正文解释与可执行实现共享同一来源。

\subsection{顺序栈：一份 size 同时分隔有效区与空闲区}
\lstinputlisting[
  language=C++,firstline=17,lastline=44,firstnumber=1,numbers=left,
  caption={顺序栈的空满、压栈、出栈与不变量}
]{30_408/10_数据结构/02_栈队列与受限访问/code/ds02_restricted_access.hpp}

代码把有效区定义为 $[0,size)$，因此压栈写 \code{data\_[size\_]} 后递增，出栈递减后读取。满/空检查在访问数组之前执行，容量为 0 时也只会进入上溢分支，不会解引用空存储。

\subsection{循环队列：每次操作只改变一个端点}
\lstinputlisting[
  language=C++,firstline=128,lastline=167,firstnumber=1,numbers=left,
  caption={循环队列的入队、出队与环形不变量}
]{30_408/10_数据结构/02_栈队列与受限访问/code/ds02_restricted_access.hpp}

入队只写 $rear$ 所指的下一空位并推进 $rear$；出队只读 $front$ 所指的首元素并推进 $front$。两者分别增减 $size$，共同维护 $rear=(front+size)\bmod capacity$。\code{next} 对零容量单独返回 0，但正常操作在此前已由空满条件拦截。

\subsection{链队列：删除唯一节点时修复两个端点}
\lstinputlisting[
  language=C++,firstline=196,lastline=240,firstnumber=1,numbers=left,
  caption={链队列的首尾迁移与不变量}
]{30_408/10_数据结构/02_栈队列与受限访问/code/ds02_restricted_access.hpp}

一般出队只移动 \code{front\_}；当旧队列只有一个节点时，移动后 \code{front\_} 为空，代码同步把 \code{rear\_} 置空。\code{invariant()} 再检查空状态等价、可达节点数、尾节点身份和尾后继为空。

\subsection{双端队列：前端与后端的顺序不对称}
\lstinputlisting[
  language=C++,firstline=269,lastline=317,firstnumber=1,numbers=left,
  caption={循环双端队列的四种端点更新}
]{30_408/10_数据结构/02_栈队列与受限访问/code/ds02_restricted_access.hpp}

\code{front\_} 始终指向首元素，\code{rear\_} 始终指向尾后下一空位。因此前插是“先退再写”，后插是“先写再进”；前删是“先读再进”，后删是“先退再读”。四个操作不是可随意对称替换的口诀，而是由端点语义生成。

\subsection{边界测试：逼状态跨过零、满与数组末端}
\lstinputlisting[
  language=C++,firstline=48,lastline=71,firstnumber=1,numbers=left,
  caption={循环队列满状态、回绕与 FIFO 断言}
]{30_408/10_数据结构/02_栈队列与受限访问/code/ds02_restricted_access_test.cpp}

测试先填满容量为 3 的队列，验证第四次入队上溢；再出队一个元素并入队 40，迫使 $rear$ 回绕，最后检查退出仍为 $10,20,30,40$ 的 FIFO 次序。测试集还覆盖零容量、空栈/空队列下溢、链队列删至空后复用，以及双端队列两端交错操作。

\begin{lstlisting}[language=bash]
clang++ -std=c++17 -Wall -Wextra -Werror -pedantic \
  code/ds02_restricted_access_test.cpp -o /tmp/ds02_test
/tmp/ds02_test
\end{lstlisting}

\section{复杂度必须绑定表示与输入契约}
\begin{center}\small
\begin{tabularx}{\linewidth}{L{2.4cm}C{1.8cm}C{1.8cm}YY}
\toprule
操作 & 顺序表示 & 链式表示 & 成本来源 & 边界条件\\
\midrule
栈 push/pop & $O(1)$ & $O(1)$ & 尾下标或头指针局部更新 & 定容上溢/动态分配\\
队列 enqueue/dequeue & $O(1)$ & $O(1)$ & 环形端点或首尾指针 & 正确空满判据\\
访问 top/front/back & $O(1)$ & $O(1)$ & 直接端点 & 非空\\
检查结构不变量 & $O(1)$ & $O(n)$ & 字段关系 vs 全链遍历 & 属于测试成本\\
存储 $n$ 个元素 & $O(M)$ & $O(n)$ & 定容槽位 vs 每节点指针 & $n\le M$\\
\bottomrule
\end{tabularx}
\end{center}
若顺序栈采用动态扩容，单次扩容为 $\Theta(n)$，连续压栈可讨论摊还 $O(1)$；本册实现采用固定容量，因而满时显式抛出上溢，不能悄悄套用动态数组结论。

\section{应用只调用时间秩序，不接管容器机制}
\begin{itemize}
\item \kw{递归/括号/表达式：}需要恢复最近尚未完成的上下文，因此调用 LIFO；文法和递归语义由相应专题解释。
\item \kw{层序遍历/BFS：}需要先展开较早发现的 frontier，因此调用 FIFO；“首次发现得到无权最短边数”由 DS07 的图模型证明。
\item \kw{DFS：}显式栈或调用栈都提供 LIFO frontier；访问标记、完成时序和图性质仍由 DS07 Own。
\item \kw{单调栈/单调队列：}它们在端点限制之外又增加“发现更优候选时淘汰旧候选”的不变量；该淘汰规则不是普通栈/队列定义。
\end{itemize}

\begin{methodbox}{Frontier Bridge 的接口语言}
调用者只需回答：新发现状态怎样加入？下一状态按什么策略取出？何时算已处理？普通 Stack/Queue 提供前两项的 LIFO/FIFO 策略；DS-B01 负责把这份接口桥接到树和图，DS04/DS07 再提供各自的正确性不变量。
\end{methodbox}

\section{概念边界：相似端点名不能替代定义}
\begin{center}\small
\begin{tabularx}{\linewidth}{L{1.8cm}C{0.35cm}L{1.8cm}YY}
\toprule
概念 A & $\neq$ & 概念 B & 真正判据 & 混淆后果\\
\midrule
ADT 次序 & $\neq$ & 物理表示 & 返回序列 vs 数组/节点布局 & 把实现细节当定义\\
rear 元素 & $\neq$ & rear 下标 & 本册 rear 指向尾后空位 & 读取错误槽位\\
队空 & $\neq$ & front=rear & 该等式必须结合消歧协议 & 把满队误判为空\\
上溢 & $\neq$ & 假溢出 & 真无槽位 vs 未复用前部槽位 & 错判容量状态\\
Deque 能力 & $\neq$ & Queue 契约 & 能两端操作 vs 只允许 FIFO & 应用可破坏时间秩序\\
容器不变量 & $\neq$ & 遍历正确性 & 结构合法 vs 算法输出正确 & Owner 越界\\
\bottomrule
\end{tabularx}
\end{center}

\begin{warnbox}{最常见的错误推理}
“循环队列就是下标取模”漏掉空满消歧；“front/rear 的公式可以直接套”漏掉二者究竟指向元素还是空位；“链式队列不会满”忽略分配失败和资源约束；“BFS 用了队列，所以最短路由队列证明”混淆容器顺序与图算法不变量。
\end{warnbox}

\section{做题控制协议与压缩复原}
遇到受限访问结构，按以下顺序建模：
\begin{enumerate}
\item 写操作集合：哪一端允许进入、离开和观察；
\item 写清端点字段语义，尤其 \code{front}/\code{rear} 指元素还是空位；
\item 选择空满消歧协议，不混用其他协议的公式；
\item 画一次普通状态和一次跨边界状态，逐步更新字段；
\item 验证抽象退出次序和物理不变量；
\item 最后把复杂度绑定到表示、容量策略和输入前提。
\end{enumerate}

\begin{corebox}{一页压缩}
忘掉公式后，从五问复原：\kw{允许动哪一端？这生成什么时间秩序？端点字段各指什么？空与满靠哪份额外信息区分？一次操作后哪些字段必须一起修复？}

栈用同端进出生成 LIFO；队列用异端进出生成 FIFO；双端队列提供更宽能力。循环化只解决槽位复用，空满消歧才决定状态是否可判；应用正确性必须继续交给调用该容器的 Topic。
\end{corebox}

\section{全科坐标与来源处理}
\begin{corebox}{Map Coordinate：Restriction / Order / State / Application}
DS02 接收 DS01 的顺序与链接表示，把“访问自由度”压缩为 LIFO/FIFO/Deque 契约；再经 DS-B01 把取出策略交给 DS04 的树遍历和 DS07 的图遍历。数据结构 Atlas 中“表示—操作—不变量—成本”在本册具体化为端点语义、状态字段与边界分支。
\end{corebox}

本册吸收归档总册中“减少自由度换取过程语义”“循环队列关键在空满规则而非取模口诀”以及 Frontier 统一视角。旧材料中 BFS/DFS、括号匹配、表达式和单调结构仅作为 Use 证据：其专属机制不提升为 DS02 的稳定定义。伴随实现承担代码级机制证据；测试承担边界与不变量证据。
\end{document}
